\subsection{Aufbau der Arbeit}

Die vorliegende Arbeit gliedert sich in insgesamt neun Kapitel. Nach der Einleitung wird in Kapitel~2 das methodische Vorgehen beschrieben, das zur Beantwortung der Forschungsfragen herangezogen wird. In Kapitel~3 werden die theoretischen Grundlagen zu \ac{ERP}-Systemen, \ac{D365FnSCM}, \ac{KI} sowie \ac{KI}-Agenten dargestellt, um ein grundlegendes Verständnis für die Thematik zu schaffen.

Darauf aufbauend wird in Kapitel~4 der Stand der Forschung zum Einsatz von \ac{KI} in \ac{ERP}-Systemen aufgearbeitet und hinsichtlich der damit verbundenen Potenziale und Herausforderungen eingeordnet. 

In Kapitel~5 wird der \ac{DSR}-Ansatz angewendet und eine Problemanalyse durchgeführt. Hierzu werden die durchgeführten Experteninterviews ausgewertet, um zunächst bestehende Probleme zu identifizieren und darauf aufbauend Anforderungen sowie erwartete Effizienzpotenziale beim Einsatz eines \ac{KI}-Agenten in \ac{D365FnSCM} abzuleiten.

Kapitel~6 widmet sich auf dieser Grundlage dem Entwurf und der Implementierung des entwickelten \ac{KI}-Agenten in \ac{D365FnSCM}. In Kapitel~7 wird der implementierte \ac{KI}-Agent anschließend hinsichtlich der zuvor identifizierten möglichen Effizienzpotenziale analysiert und bewertet.

Abschließend werden in Kapitel~8 die zentralen Ergebnisse der Arbeit zusammengefasst und die Limitationen diskutiert. Kapitel~9 schließt die Arbeit mit der Beantwortung der Forschungsfragen sowie einem Ausblick auf mögliche zukünftige Forschungsfelder ab.
